\section{Related Work}
\label{sec:related_work}

\para{LayerNorm as more than an optimization trick.} Early transformer work largely treated LayerNorm as part of the recipe needed for stable training, and \citet{xiong2020layernorm} gave one of the clearest analyses of how normalization placement affects gradients and warm-up behavior. More recent work argues that LayerNorm also changes model expressivity. \citet{brody2023expressivity} decompose its geometric role and show that it affects which keys become selectable by attention. \citet{zhu2025transformers} extend the challenge further by showing that strong transformer families can work without standard normalization layers. These papers motivate our hypothesis, but they do not test whether the geometric role of LayerNorm matters for creative generation or semantic association.

\para{Degeneration and diversity in open-ended generation.} Work on neural text degeneration shows that high-probability decoding strategies can produce bland and repetitive outputs even when the underlying model is strong \citep{holtzman2020curious}. More recently, \noveltybench demonstrates that modern language models often generate far less distributional diversity than humans, even when they remain useful or coherent \citep{zhang2025noveltybench}. This line of work frames the output-side problem clearly, but it usually studies decoding, prompting, or model scale rather than a targeted normalization ablation.

\para{Creativity and remote association benchmarks.} Creativity studies have started to measure language models with cognitive tasks rather than only subjective judgments. \citet{chen2023creativity} use the Divergent Association Task to show that stochastic decoding can raise semantic-distance scores for some models, though often with a stability trade-off. \citet{stevenson2022gpt3} evaluate GPT-3 on the Alternative Uses Test and find that humans still outperform the model on overall creative quality. \citet{haase2025creativity} emphasize that repeated prompting reveals large intra-model variability, which motivates our multi-sample protocol. In parallel, \citet{naeini2023ocw} introduce the Only Connect Wall benchmark as a distractor-rich association task that exposes brittle grouping behavior in language models. Our study builds on these benchmarks but asks a different question: whether LayerNorm changes the geometry that supports or constrains such behaviors.

\para{Positioning our work.} Unlike prior work, we do not compare unrelated model families or alternate only the decoder. We compare matched GPT-2 variants that differ in LayerNorm status, measure both external creativity and internal association geometry, and test whether the latter predicts the former. This is the missing bridge between normalization mechanics and creative-generation outcomes.
